% GENERATED FILE. Edit canonical lessons, then run npm run generate:books.
% canonical-source-hash: fnv1a64:0b143c7b8ffea46b
% canonical-lessons: LA-C550-litteratus, LA-C550-bibliopola, LA-C550-lector, LA-C550-perlego, LA-C550-evolvo

\chapter{Readers and booksellers}
\label{ch:la-readers-and-booksellers}
\hlchaptermodality{550}

\begin{chapteropening}
\textbf{By the end of this chapter:} \emph{I can talk about readers, booksellers and reading.}

Complete the last lesson of chapter 550: I can talk about readers, booksellers and reading.
\end{chapteropening}

\section[litterātus, litterāta, litterātum]{litterātus, litterāta, litterātum — learned, well-read}

\label{lesson:LA-C550-litteratus}



\begin{quote}
Before the new one: say the Latin for a public reading, then the Latin for a listener.
\end{quote}

\subsection*{You'll want to know: litterātus, litterāta, litterātum}
\textbf{litterātus, litterāta, litterātum} — \textquotedblleft{}learned, well-read\textquotedblright{}.

\begin{quote}
Vir litterātus multōs librōs habet.
\end{quote}

\begin{quote}
\emph{The learned man has many books.}
\end{quote}

From \emph{litterae} — letters, learning. English \emph{literate}.

\begin{cousinweb}
\textbf{litterātus} is built on \emph{littera} \textquotedblleft{}a letter\textquotedblright{} (Lewis \& Short; EtymDB): lettered, learned. English \emph{literate} comes from it.
\end{cousinweb}

\begin{grammarlens}[title={What you've built}]
One more word for telling, explaining and describing at B1.
\end{grammarlens}

\subsection*{Your turn}
\begin{itemize}
\raggedright
  \item \emph{Say it:} \emph{litterātus, litterāta, litterātum}
  \item \emph{Say it:} \emph{litterātus, litterāta, litterātum}, once more
  \item \emph{Say it:} a sentence of your own with \emph{litterātus, litterāta, litterātum}
  \item \emph{Recall:} say the Latin for a public reading, then the Latin for a listener, then say \emph{litterātus, litterāta, litterātum} again
\end{itemize}

\begin{tcolorbox}[breakable,colback=teal!4,colframe=teal!35!black,title={Before you move on}]
What does litterātus, litterāta, litterātum mean? (\textquotedblleft{}Learned, well-read\textquotedblright{}.) Say it once more.
\end{tcolorbox}
\section[bibliopōla, bibliopōlae]{bibliopōla, bibliopōlae — a bookseller}

\label{lesson:LA-C550-bibliopola}



\begin{quote}
Before the new one: say the Latin for a listener, then the Latin for learned.
\end{quote}

\subsection*{You'll want to know: bibliopōla, bibliopōlae}
\textbf{bibliopōla, bibliopōlae} — \textquotedblleft{}a bookseller\textquotedblright{}.

\begin{quote}
Bibliopōla novōs librōs in tabernā vendit.
\end{quote}

\begin{quote}
\emph{The bookseller sells new books in his shop.}
\end{quote}

Masculine, first declension, from Greek \emph{biblion} 'book' + \emph{pōlein} 'to sell'.

\begin{cousinweb}
\textbf{bibliopōla} was borrowed from Greek \emph{bibliopoles}, built on \emph{biblion} \textquotedblleft{}book\textquotedblright{} and \emph{poles} \textquotedblleft{}seller\textquotedblright{} (EtymDB; Lewis \& Short).
\end{cousinweb}

\begin{grammarlens}[title={What you've built}]
One more word for telling, explaining and describing at B1.
\end{grammarlens}

\subsection*{Your turn}
\begin{itemize}
\raggedright
  \item \emph{Say it:} \emph{bibliopōla, bibliopōlae}
  \item \emph{Say it:} \emph{bibliopōla, bibliopōlae}, once more
  \item \emph{Say it:} a sentence of your own with \emph{bibliopōla, bibliopōlae}
  \item \emph{Recall:} say the Latin for a listener, then the Latin for learned, then say \emph{bibliopōla, bibliopōlae} again
\end{itemize}

\begin{tcolorbox}[breakable,colback=teal!4,colframe=teal!35!black,title={Before you move on}]
What does bibliopōla, bibliopōlae mean? (\textquotedblleft{}A bookseller\textquotedblright{}.) Say it once more.
\end{tcolorbox}
\section[lēctor, lēctōris]{lēctor, lēctōris — a reader}

\label{lesson:LA-C550-lector}



\begin{quote}
Before the new one: say the Latin for learned, then the Latin for a bookseller.
\end{quote}

\subsection*{You'll want to know: lēctor, lēctōris}
\textbf{lēctor, lēctōris} — \textquotedblleft{}a reader\textquotedblright{}.

\begin{quote}
Lēctor dominō librum legit.
\end{quote}

\begin{quote}
\emph{The reader reads a book to his master.}
\end{quote}

Masculine, third declension, from \emph{legō}. Rich houses kept a slave to read aloud.

\begin{cousinweb}
\textbf{lēctor} is built on \emph{legō} \textquotedblleft{}to read\textquotedblright{} with \emph{-tor} (Lewis \& Short; EtymDB). English \emph{lecture} comes from the related \emph{lēctūra}.
\end{cousinweb}

\begin{grammarlens}[title={What you've built}]
One more word for telling, explaining and describing at B1.
\end{grammarlens}

\subsection*{Your turn}
\begin{itemize}
\raggedright
  \item \emph{Say it:} \emph{lēctor, lēctōris}
  \item \emph{Say it:} \emph{lēctor, lēctōris}, once more
  \item \emph{Say it:} a sentence of your own with \emph{lēctor, lēctōris}
  \item \emph{Recall:} say the Latin for learned, then the Latin for a bookseller, then say \emph{lēctor, lēctōris} again
\end{itemize}

\begin{tcolorbox}[breakable,colback=teal!4,colframe=teal!35!black,title={Before you move on}]
What does lēctor, lēctōris mean? (\textquotedblleft{}A reader\textquotedblright{}.) Say it once more.
\end{tcolorbox}
\section[perlegō, perlegere]{perlegō, perlegere — to read through}

\label{lesson:LA-C550-perlego}



\begin{quote}
Before the new one: say the Latin for a bookseller, then the Latin for a reader.
\end{quote}

\subsection*{You'll want to know: perlegō, perlegere}
\textbf{perlegō, perlegere} — \textquotedblleft{}to read through\textquotedblright{}.

\begin{quote}
Epistulam tuam ter perlēgī.
\end{quote}

\begin{quote}
\emph{I have read your letter through three times.}
\end{quote}

\emph{perlegō, perlegere, perlēgī, perlēctum}: \emph{per} 'through' + \emph{legō}.

\begin{cousinweb}
\textbf{perlegō} is built on \emph{per-} \textquotedblleft{}through\textquotedblright{} and \emph{legō} \textquotedblleft{}to read\textquotedblright{} (EtymDB; Lewis \& Short).
\end{cousinweb}

\begin{grammarlens}[title={What you've built}]
One more word for telling, explaining and describing at B1.
\end{grammarlens}

\subsection*{Your turn}
\begin{itemize}
\raggedright
  \item \emph{Say it:} \emph{perlegō, perlegere}
  \item \emph{Say it:} \emph{perlegō, perlegere}, once more
  \item \emph{Say it:} a sentence of your own with \emph{perlegō, perlegere}
  \item \emph{Recall:} say the Latin for a bookseller, then the Latin for a reader, then say \emph{perlegō, perlegere} again
\end{itemize}

\begin{tcolorbox}[breakable,colback=teal!4,colframe=teal!35!black,title={Before you move on}]
What does perlegō, perlegere mean? (\textquotedblleft{}To read through\textquotedblright{}.) Say it once more.
\end{tcolorbox}
\section[ēvolvō, ēvolvere]{ēvolvō, ēvolvere — to unroll; to read (a book)}

\label{lesson:LA-C550-evolvo}



\begin{quote}
Before the new one: say the Latin for a reader, then the Latin for to read through.
\end{quote}

\subsection*{You'll want to know: ēvolvō, ēvolvere}
\textbf{ēvolvō, ēvolvere} — \textquotedblleft{}to unroll; to read (a book)\textquotedblright{}.

\begin{quote}
Discipulus librum ēvolvit et legit.
\end{quote}

\begin{quote}
\emph{The pupil unrolls the book and reads.}
\end{quote}

\emph{ēvolvō, ēvolvere, ēvolvī, ēvolūtum}. Books were rolls, so to read was to unroll. English \emph{evolve}.

\begin{cousinweb}
\textbf{ēvolvō} is built on \emph{ē-} \textquotedblleft{}out\textquotedblright{} and \emph{volvō} \textquotedblleft{}to roll\textquotedblright{} (EtymDB; Lewis \& Short): to unroll a book. English \emph{evolve} comes from it.
\end{cousinweb}

\begin{grammarlens}[title={What you've built}]
One more word for telling, explaining and describing at B1.
\end{grammarlens}

\subsection*{Your turn}
\begin{itemize}
\raggedright
  \item \emph{Say it:} \emph{ēvolvō, ēvolvere}
  \item \emph{Say it:} \emph{ēvolvō, ēvolvere}, once more
  \item \emph{Say it:} a sentence of your own with \emph{ēvolvō, ēvolvere}
  \item \emph{Recall:} say the Latin for a reader, then the Latin for to read through, then say \emph{ēvolvō, ēvolvere} again
\end{itemize}

\begin{tcolorbox}[breakable,colback=teal!4,colframe=teal!35!black,title={Before you move on}]
What does ēvolvō, ēvolvere mean? (\textquotedblleft{}To unroll; to read (a book)\textquotedblright{}.) Say it once more.
\end{tcolorbox}
